\section{Introduction}
\label{sec:intro}

Ionic liquids (ILs) are salts that are liquid below about 100\,\textdegree C. Their negligible vapour pressure, thermal stability and the freedom to pair many cations with many anions make them candidate solvents for \co{} capture \cite{blanchard1999green, brennecke2010ionic, zhang2012carbon, aghaie2018systematic}. In common imidazolium ILs, \co{} dissolves mainly by physical absorption, and the anion has the largest effect on solubility \cite{cadena2004why, anthony2005anion}. Functionalised ILs can also bind \co{} chemically \cite{bates2002co2, gurkan2010equimolar}. The number of possible cation--anion pairs is very large, while each solubility isotherm requires dedicated high-pressure measurements. A model that predicts solubility for ILs that have not been measured would let experiments focus on the most promising candidates.

Data-driven models have been built for this task, from early neural networks \cite{eslamimanesh2011artificial} to tree ensembles and deep networks trained on several thousand points \cite{venkatraman2017predicting, song2020prediction, ali2024prediction}. Most use group-contribution or COSMO-derived inputs. This work instead describes each ion with a small set of descriptors computed by density functional theory (DFT), so that any ion whose structure is known can be described without experimental input. A further motivation is local: Indonesia is a major producer of palm and coconut oil, whose fatty acids (oleic and lauric acid) could supply the anions of bio-based ILs.

The contributions of this work are:
\begin{enumerate}
    \item an experimental dataset of 11,256 \co{} solubility points for 113 ILs, and a rebuilt version with full source references (Section~\ref{sec:data});
    \item a pipeline from ion names to SMILES, 3D structures and r2SCAN-3c calculations, giving five DFT descriptors per ion (Section~\ref{sec:comp});
    \item a comparison of four regressors and a feature-importance analysis (Sections~\ref{sec:models} and~\ref{sec:results});
    \item a first screening of imidazolium ILs and of two bio-based ILs (Section~\ref{sec:screening}).
\end{enumerate}
Section~\ref{sec:issues} reports the limitations found when the workflow was checked against the raw calculation outputs, and Section~\ref{sec:summary} concludes.
